% Part: proof-theory
% Chapter: natural-deduction
% Section: sequents

\documentclass[../../../include/open-logic-section]{subfiles}

\begin{document}

\olfileid{pt}{nat}{seq}

\olsection{带相继式的自然演绎：系统 \Log{N2c} 与~\Log{N2i}}

\Log{G1c} 或 \Log{G1i} 中的一个推导~$\delta$ 表明，它的结束公式~$!A$ 可以从其开放假设推出。
如果 $\Gamma$ 是由所有满足如下条件的公式~$!B$ 构成的集合：$!B^x$ 在~$\delta$ 中作为开放假设出现，
那么我们可以将此事记作 $\delta \Proves \Gamma \Sequent !A$。
这表明我们也可以为相继式构造一个自然演绎系统。
在这样的系统中，每个相继式都承载了右边公式依赖于哪些开放假设的信息，
也就是在以该相继式结束的子推导中，哪些假设仍是开放的。
这就是\emph{带相继式的自然演绎}，或称“相继式风格的”自然演绎，
我们将得到的系统称为 \Log{N2c} 和~\Log{N2i}。

为了让 \Log{N1} 和 \Log{N2} 之间的对应更紧密，
我们令左边的集合~$\Gamma$ 由带标记的公式~$x:!B$ 构成。
由于 \Log{N1} 的规则只作用于前提中的公式（即直接子推导的结束公式），
所以 \Log{N2} 的规则也只作用于相继式的后件。
因此，我们可以简单地将相继式左边带标记的公式称为\emph{上下文}。

\Log{N2} 的证明不以假设作为起点，而是以形如
\[x:!A \Sequent !A.\]
的公理作为最顶层的相继式。
其规则与 \Log{N1} 的规则完全对应，具体见\olref{tab:N2}。

\subfile{rules-N2}

由于我们必须允许（正如我们在 \Log{N1c} 和 \Log{N1i} 中所做的那样）规则 \Intro{\lif}、\Intro{\lnot}、\Elim{\lor} 和 \Elim{\lexists} 可以但并不必须实际解除假设，因此我们也允许在 \Log{N2c} 和~\Log{N2i} 中正确地应用相应的规则，即使相应的上下文公式 $!A$、$!B$ 和 $!A(a)$ 并未出现在相关前提的上下文中。

\begin{prop}
如果 $\delta$ 是 \Log{N1c} 或 \Log{N1i} 中~$!A$ 的一个推导，那么就存在 \Log{N2c}（或 \Log{N2i}）中 $\Gamma \Sequent !A$ 的一个推导~$\delta'$，其中 $\Gamma$ 是由所有满足 $!B^x$ 是~$\delta$ 的开放假设的~$x:!B$ 构成的集合。
\end{prop}

\begin{proof}
  对 $n=\pheight{\delta}$ 进行归纳。若 $n=0$，则 $\delta$ 是一个单独的开放假设~$!A^x$。那么 $\Gamma = \{x:!A\}$ 且 $\Gamma \Sequent !A$ 是 \Log{N2c} 或~\Log{N2i} 的一条公理。

  若 $n>0$，我们根据~$\delta$ 的最后一步推理区分情况。我们只讨论一种情况；其余留作练习。

  假设最后一步推理是 $\Intro{\lif}$。那么 $\delta$ 以
  \[
  \AxiomC{$\Discharge{!B}{x}$}
  \RightLabel{$\delta_1$}
  \DeduceC{$!C$}
  \DischargeRule{\Intro{\lif}}{x}
  \UnaryInfC{$!B \lif !C$}
  \DisplayProof
  \]
  结束。根据归纳假设，存在 \Log{N2c}（或 \Log{N2i}）中的一个推导~$\delta_1'$，其结论为 $\Gamma' \Sequent !C$，其中 $\Gamma'$ 是~$\delta_1$ 中开放的带标记假设。如果 $!B^x$ 是 $\delta_1$ 的一个开放假设，那么它在 \Intro{\lif} 推理处被解除，并且~$\delta$ 的开放假设是 $\Gamma =
  \Gamma' \setminus \{x:!B\}$。换句话说，$\delta_1'$ 是
  $x:!B, \Gamma \Sequent !C$ 的一个推导，并且我们可以将 $\delta'$ 取为
  \[
  \AxiomC{}
  \RightLabel{$\delta_1'$}
  \Deduce$x:B, \Gamma \fCenter !C$
  \RightLabel{\Intro{\lif}}
  \UnaryInf$\Gamma \fCenter !B \lif !C$
  \DisplayProof
  \]
  如果 $!B^x$ 不是~$\delta_1$ 的开放假设，那么 \Intro{\lif} 并未解除任何假设，并且 $\delta$ 与 $\delta_1$ 的开放假设相同。由于在 \Log{N2c} 和~\Log{N2i} 中，并不要求上下文公式 $x:!B$ 出现在 \Intro{\lif} 的前提中，我们可以将 $\delta'$ 取为
  \[
  \AxiomC{}
  \RightLabel{$\delta_1'$}
  \Deduce$\Gamma \fCenter !C$
  \RightLabel{\Intro{\lif}}
  \UnaryInf$\Gamma \fCenter !B \lif !C$
  \DisplayProof
  \]

  $\delta$ 以其他规则结束的情况是类似的。
\end{proof}

\begin{prop}
如果 $\delta$ 是 \Log{N2c} 或 \Log{N2i} 中 $\Gamma \Sequent !A$ 的一个推导，那么存在 \Log{N1c}（或 \Log{N1i}）中的一个推导~$\delta'$，其结束公式为~$!A$，且对于每个 $x:!B \in \Gamma$，$!B^x$ 都是~$\delta'$ 的开放假设。
\end{prop}

\begin{proof}
  我们再次对~$\delta$ 的高度~$n$ 进行归纳。若
  $n=0$，则 $\delta$ 是一条公理 $x:!A \Sequent !A$。
  推导~$\delta'$ 就是假设~$!A^x$。

  若 $n>0$，我们根据~$\delta$ 的最后一步推理区分情况。例如，如果该推理是~$\Intro{\lif}$，
  $\delta$~以
  \[
  \bottomAlignProof
  \AxiomC{}
  \RightLabel{$\delta_1$}
  \Deduce$x:!B, \Gamma \fCenter !C$
  \RightLabel{\Intro{\lif}}
  \UnaryInf$\Gamma \fCenter !B \lif !C$
  \DisplayProof
  \quad\text{ 或 }\quad
  \bottomAlignProof
  \AxiomC{}
  \RightLabel{$\delta_1$}
  \Deduce$\Gamma \fCenter !C$
  \RightLabel{\Intro{\lif}}
  \UnaryInf$\Gamma \fCenter !B \lif !C$
  \DisplayProof
  \]
  结束。根据归纳假设，存在 \Log{N1c}（或 \Log{N1i}）中~$!C$ 的一个推导。在第一种情况下，$!B^x$ 是~$\delta'$ 的一个开放假设；在第二种情况下则不是。我们可以将推导~$\delta'$ 相应地取为
  \[
  \bottomAlignProof
  \AxiomC{$\Discharge{!B}{x}$}
  \RightLabel{$\delta_1$}
  \DeduceC{$!C$}
  \DischargeRule{\Intro{\lif}}{x}
  \UnaryInfC{$!B \lif !C$}
  \DisplayProof
  \quad\text{ 或 }\quad
  \bottomAlignProof
  \AxiomC{}
  \RightLabel{$\delta_1$}
  \DeduceC{$!C$}
  \RightLabel{\Intro{\lif}}
  \UnaryInfC{$!B \lif !C$}
  \DisplayProof
  \]
\end{proof}

\end{document}